\section{Conclusion}
\label{sec:conclusion}

In this Research Preview, we studied how chunking strategy and context headers affect requirement-to-code evidence retrieval over 40 repositories. 
After reranking, headers improved AST chunks ($\Delta = +0.112$, $p_{\mathrm{Holm}} = 0.0031$); without headers, line windows outperformed AST chunks ($\Delta = -0.126$, $p_{\mathrm{Holm}} = 0.0001$); and we found no evidence that AST chunks with headers differ from line windows without headers ($\Delta = -0.014$, $p_{\mathrm{Holm}} = 0.85$). 
Line windows with headers scored highest (NDCG@10 = 0.499 vs. 0.457), though this difference was exploratory and not significant. 
Because 35.7\% of AST-arm chunks are line-window fallbacks, these results concern an AST-first arm with fallback, not pure AST chunking. 
To invite discussion, we pose three questions: 
(1)~\textit{Is structural chunking worth its parsing complexity once context headers are available?} 
(2)~\textit{How should code-RAG benchmarks balance syntactic boundary integrity against lexical content?} 
(3)~\textit{How should evidence retrieval for requirements be used responsibly when the requirements describe people's skills?}

\section*{Data Availability}
The package includes chunk text from 40 public repositories (19 MIT, 4 Apache-2.0, 1 AGPL-3.0, 16 without a stated license); copyright remains with the original authors, and the manifest lists repository URLs and commit SHAs. 
The package is accessible at: \url{https://anonymous.4open.science/r/ast-rag-replication-2027}.

\section*{Acknowledgments}
Generative AI tools (Google Antigravity, including Anthropic Claude and Google Gemini models) were used for drafting and editing text, and for generating analysis scripts, tables, and figures. 
All AI-generated text and code were reviewed, verified, and edited by the authors, who take full responsibility for the content. 
Relevance labels were produced by human annotators without AI assistance. 
An LLM was also used to extract queries from job descriptions, as described in Section~\ref{sec:approach}.
